
\edps\ is an environment designed to execute the recipes of an instrument pipeline according to a series of instructions. 
The main concepts in \edps\ are:

\begin{itemize}
  
\item \textbf{Workflow} and \textbf{reduction cascades}. A workflow is
  a series of instructions designed to reduce data with an instrument
  pipeline in potentially multiple ways, by carrying on a sequence of
  tasks.  Each workflow can define multiple reduction cascades,
  depending on the scientific needs.  For example, the same workflow
  can be used to process data following different strategies that
  trigger different reduction steps (e.g. in one strategy flux
  calibration can be omitted) or different end-points (e.g., combine
  different science exposures, or stop after the reduction of
  individual exposures without combining them). Each of these
  "strategies" defines a "reduction cascade".


\item \textbf{Task}, \textbf{jobs}, and \textbf{recipes}. A task is an
  element in the workflow that performs a given step of the data
  reduction cascade. Tasks are often associated to a recipe of the
  underlying instrument pipeline.  A jobs is a work unit in a
  processing environment, that runs a recipe on a set of input data
  with a set of recipe parameters.  A single task can generate several
  jobs: for example, a "bias" task, can generate multiple jobs, each
  of the running the bias recipe on a different set of input files.


\item \textbf{Dataset}. A dataset is a collection of files, that are
  needed to perform the data reduction as specified by the
  workflow. It consists, for example, of one or more science files
  plus the calibrations needed to process them. In \edps,  datasets
  have a hierarchical structure, which highlights the connections
  between the various files and tasks (e.g., task A is an input to
  task B).


\item \textbf{Target} and \textbf{Target category}. The "target", or
  the "target task" is the end point of the reduction cascade. When
  specifying a target, \edps\ will process all and only the files needed
  to execute it. For example, if my target is "science", and the
  science files need the bias files, \edps\ will process only the biases
  that have been selected to process those science files; then it
  processes the science using the product of the bias reduction.
  However, if my target is bias, then \edps\ will process all and only
  the bias files, regardless they are not used by any science.  In
  this case, \edps\ does not process the science, as it has already reached
  the end reduction point (e.g., process all biases).  The
  "Target category" is a group of targets that have similar
  purposes. For example, the target category "science", includes all
  the tasks that deliver final scientific products, the target
  category "qc1calib" includes all and only the tasks that processes
  calibrations (e.g., bias, flat fields, standard stars).

\end{itemize}
